\chapter{LINE Messaging API 簡介}

LINE Messaging API 是 LINE 平台提供的一套用於建立自動化對話、通知與互動功能的開發介面。它讓開發者可以設計智慧型機器人、客服系統、通知服務與自動回覆流程，並與使用者在 LINE 上建立雙向溝通。就像許多入口網站與行動服務都需透過 API 與後端系統整合，LINE Messaging API 也為企業與個人開發者提供穩定、可擴充的訊息通路。

在實務層面，LINE Messaging API 可應用於客服排班、會員通知、活動促銷、問卷回覆、訂單追蹤與內部溝通等情境。若公司希望透過聊天平台提供服務，而不必另行開發完整 APP，這種 API 結合 Webhook 與後端程式的方式相當實用。其優點在於整合成本較低、使用者不需額外下載新工具，且平台上的訊息互動模式容易被接受。

\section{LINE Messaging API 是什麼？}

LINE Messaging API 其實是一組基於 HTTP 與 Webhook 的開發標準。開發者先建立一個 Channel，取得必要的認證資訊，接著讓後端伺服器接收 LINE 平台發送的事件，例如使用者送出文字、貼圖、圖片或定位資訊。當事件發生時，伺服器可根據條件判斷並回覆文字、圖片、卡片式訊息，甚至進行推播。

這種開發方式的核心在於「事件驅動」。也就是說，當使用者在 LINE 上輸入內容、點擊按鈕或發送圖片時，LINE 會將事件資訊轉送到 Webhook URL，開發者就能接收這些事件，並做相對應的業務邏輯處理。這使得 LINE 不只是聊天軟體，而是可與企業系統整合的溝通平臺。

\section{常見應用場景}

LINE Messaging API 的應用並不限於聊天機器人，以下是幾種比較典型的使用方式：

\begin{itemize}[leftmargin=2em]
    \item 自動回覆客服：例如會員詢問營業時間、訂單狀態或退貨流程。
    \item 推播通知：例如活動提醒、發票通知、停車場剩餘車位提醒。
    \item 互動式問卷：使用者在 LINE 中輸入選項，系統自動記錄並分析資料。
    \item 校園或公司內部管理：通知公告、預約服務、活動報名與審核。
    \item 電商與訂單管理：例如購物車提醒、配送狀態更新與售後回復。
\end{itemize}

可以說，LINE Messaging API 不僅是「聊天」功能，而是把過去需要使用者主動開啟應用程式的操作，轉變成以聊天平台為入口的服務體驗。這也是為什麼它在行銷、客服、教育與企業內部通訊中都具有相當大的發展潛力。

\begin{figure}[H]
\centering
\fbox{\parbox[c][2.8cm][c]{12cm}{\centering LINE Developer Console 示意圖\\（可在實際平台中替換成截圖）}}
\caption{LINE Developer Console 開發者入口}
\end{figure}

\begin{table}[H]
\centering
\caption{LINE Messaging API 常見功能比較}
\begin{tabular}{p{4.5cm}p{7.5cm}}
\toprule
功能 & 說明 \\
\midrule
訊息推播 & 後端主動發送通知訊息給使用者 \\
自動回覆 & 接收到使用者訊息時，依規則回覆內容 \\
Rich Menu & 客製化主選單，提升使用者操作效率 \\
Flex Message & 高度自訂的卡片式訊息 \\
Webhook & 由 LINE 平台將事件回調到指定伺服器 \\
\bottomrule
\end{tabular}
\end{table}

\chapter{LINE Developers 註冊與 Channel 建立}

若要使用 LINE Messaging API，第一步是建立 LINE Developers 帳號，並在 Developer Console 中建立 Provider 與 Messaging API Channel。整個流程可視為「申請開發者身分」與「建立機器人通道」的雙重動作。完成後，系統將提供必要的 Channel Secret、Channel Access Token、Webhook URL 等資訊。

\section{註冊步驟}

\begin{enumerate}[leftmargin=2em]
    \item 進入 LINE Developers 官方網站，使用 LINE 帳號登入。
    \item 點選「Create a provider」建立 Provider，這相當於專案或組織的容器。
    \item 在 Provider 內建立「Messaging API」Channel。
    \item 填寫 Channel 名稱、說明、功能類型與相關資訊。
    \item 完成建立後，即可取得識別碼、Secret 與 Access Token。
\end{enumerate}

這些步驟看似簡單，但對正式專案而言，準備完整的專案資訊非常重要。因為 Channel 是後續 API、Webhook 與機器人行為的核心，而 Provider 則是用來管理一組相關 Channel 的單位。若您未來要開發多個業務場景，例如客服機器人與訂閱通知，也可以在同一個 Provider 下建立不同 Channel。

\section{Channel 建立的實務要點}

在建立 Channel 時，建議先確認以下內容：

\begin{itemize}[leftmargin=2em]
    \item Channel 名稱是否容易辨識，方便未來管理多個機器人。
    \item 是否需要開啟 Webhook 功能，進行事件回調。
    \item 是否要支援 Rich Menu、Quick Reply 或 Flex Message 等高互動訊息格式。
    \item 是否需要設定「自動回覆」或「訊息推播」權限。
\end{itemize}

完成 Channel 建立後，開發者通常會得到一個專用的 Channel ID、Channel Secret 以及 Access Token。這些資訊是後續與 LINE 伺服器整合的核心密碼，務必要妥善保管，避免直接放在公開程式碼中。

\chapter{Channel Secret、Access Token 與安全設定}

在 LINE Messaging API 中，最重要的安全資源大致可分為兩類：Channel Secret 與 Access Token。前者相當於機器人的基本識別資訊，後者則更像另一個可供程式端存取 LINE API 的憑證。這些資料的取得與保護方式直接影響整個聊天機器人的安全。

\section{Channel Secret}

Channel Secret 是 LINE 控制台中隱藏的識別金鑰，用於確認請求來源是否為合法的 LINE 伺服器或官方堆疊。通常在建立 Channel 後，開發者可在「Basic settings」中查閱。由於它充當重要的驗證資訊，若在程式碼或公開環境中直接明碼儲存，將使系統容易遭受串接風險。

\section{Access Token}

Access Token 則是讓後端程式「呼叫 LINE API」的授權憑證，例如發送訊息、建立 Rich Menu 或取得資料。它通常分為長期令牌與短期令牌兩種設計，開發者可依實際功能需求來管理和更新。如果未妥善保護，攻擊者可能透過這些令牌冒充機器人發送通知，造成資安與品牌風險。

\begin{lstlisting}[language=Python, caption={以環境變數方式儲存 LINE 認證資訊}, label={lst:secure}]
import os

CHANNEL_SECRET = os.getenv("LINE_CHANNEL_SECRET")
CHANNEL_ACCESS_TOKEN = os.getenv("LINE_CHANNEL_ACCESS_TOKEN")

if not CHANNEL_SECRET or not CHANNEL_ACCESS_TOKEN:
    raise RuntimeError("Missing LINE credentials")
\end{lstlisting}

\section{安全建議}

\begin{enumerate}[leftmargin=2em]
    \item 將憑證放入環境變數，而不是直接寫死於程式碼中。
    \item 不要將 Token 或 Secret 提交到公開 Git 倉庫。
    \item 在正式環境中使用 HTTPS 與有效的伺服器憑證。
    \item 對 Webhook 進行簽章驗證，避免未驗證請求進入伺服器。
    \item 定期輪換 Access Token，降低被竊取後的影響範圍。
\end{enumerate}

這些安全準則不只是基本最佳實務，也是開發 LINE 機器人時不可忽略的核心。

\chapter{Webhook 設定與訊息接收}

Webhook 是 LINE Messaging API 的關鍵能力之一。它允許 LINE 平台在使用者發送訊息、點擊按鈕或觸發其他事件時，將事件資訊即時回傳至指定的伺服器 URL。這使得開發者能夠在後端主動處理使用者行為，而不是只能被動等待使用者請求。

\section{Webhook 的工作方式}

通常，Webhook 設定流程如下：

\begin{enumerate}[leftmargin=2em]
    \item 開發者在 LINE Developer Console 啟用 Webhook。
    \item 伺服器提供一個公開 HTTPS URL，例如 \texttt{https://example.com/callback}。
    \item LINE 在收到使用者事件後，將 JSON 格式事件傳送到該 URL。
    \item 後端程式驗證簽章，再解析事件內容並進行相應處理。
\end{enumerate}

以客服或問答機器人為例，使用者輸入「查詢訂單」之後，LINE 將訊息事件送到後端；伺服器收到事件後，根據關鍵字進行查詢，再回覆結果。這種架構極適合自動化問答、工作通知與智慧型對話整合。

\begin{lstlisting}[language=Python, caption={Flask Webhook 接收範例}, label={lst:webhook}]
from flask import Flask, request, abort
from linebot import LineBotApi, WebhookHandler
from linebot.exceptions import InvalidSignatureError
from linebot.models import MessageEvent, TextMessage, TextSendMessage

app = Flask(__name__)
line_bot_api = LineBotApi("YOUR_CHANNEL_ACCESS_TOKEN")
handler = WebhookHandler("YOUR_CHANNEL_SECRET")

@app.route("/callback", methods=["POST"])
def callback():
    signature = request.headers.get("X-Line-Signature", "")
    body = request.get_data(as_text=True)
    try:
        handler.handle(body, signature)
    except InvalidSignatureError:
        abort(400)
    return "OK"

@handler.add(MessageEvent, message=TextMessage)
def handle_text_message(event):
    reply = f"你剛剛說了：{event.message.text}"
    line_bot_api.reply_message(event.reply_token, TextSendMessage(text=reply))

if __name__ == "__main__":
    app.run(host="0.0.0.0", port=5000)
\end{lstlisting}

\section{ngrok 測試方法}

在本地開發時，伺服器通常位於內網或執行環境中，無法直接被 LINE 平台訪問。因此許多開發者使用 \texttt{ngrok} 將本機服務公開到外網。其流程相當簡單：

\begin{itemize}[leftmargin=2em]
    \item 啟動本地 Flask 應用程式。
    \item 執行 \texttt{ngrok http 5000}。
    \item 將轉發網址填入 LINE Webhook 設定。
    \item 使用真實 LINE 帳號測試訊息互動。
\end{itemize}

這種方法非常實用，因為它能在不花費額外伺服器成本的前提下，快速完成功能驗證與除錯。

\chapter{LINE 訊息推播與自動回覆實作}

在完成 Channel 認證與 Webhook 設定之後，開發者即可開始實作「自動回覆」與「主動推播」兩種主要功能。前者是接收到事件後立即回應，後者則是在特定條件成立時主動發送訊息。這兩者是 LINE 機器人最常見的核心功能。

\section{接收與回覆訊息}

以下範例展示如何在收到使用者文字訊息時，回覆一段固定內容，並依據使用者輸入字串做判斷：

\begin{lstlisting}[language=Python, caption={接收使用者文字訊息並回覆範例}, label={lst:reply}]
from linebot.models import MessageEvent, TextMessage, TextSendMessage

@handler.add(MessageEvent, message=TextMessage)
def handle_message(event):
    msg = event.message.text

    if msg == "你好":
        reply_text = "你好！我是 LINE Bot，請問有什麼可以幫你？"
    elif msg in ["時間", "現在時間"]:
        import datetime
        now = datetime.datetime.now().strftime("%Y-%m-%d %H:%M:%S")
        reply_text = f"現在時間是 {now}"
    else:
        reply_text = f"你剛剛說的是：{msg}"

    line_bot_api.reply_message(
        event.reply_token,
        TextSendMessage(text=reply_text)
    )
\end{lstlisting}

這種程式邏輯很適合用在問答機器人、客服導引與簡單自動化流程。若想實現更複雜的交互，例如關鍵字分類、會員查詢、訂單狀態查詢或客服排班，則需在後端整合資料庫與業務邏輯。

\section{推播訊息與進階功能}

除了回覆訊息外，LINE 也支援進行 Push Message。這表示系統可以在不等待使用者發送訊息的情況下，主動發送通知。例如：

\begin{itemize}[leftmargin=2em]
    \item 活動開始提醒
    \item 訂單狀態更新
    \item 系統維護公告
    \item 會員生日祝福
\end{itemize}

進一步說，LINE 也支援 Rich Menu、Quick Reply、Template Message 與 Flex Message 等高互動內容。這些介面可提升使用者體驗，讓機器人不只是純文字回覆，而是可呈現按鈕、卡片、選單與豐富內容。對於商業應用而言，這種設計能明顯提升行銷與客服效率。

\section{實務建議}

\begin{enumerate}[leftmargin=2em]
    \item 先從最小可行功能開始，例如「問候語」與「關鍵字回覆」。
    \item 完成 Webhook 驗證後，再逐步加入資料庫與外部 API。
    \item 建立日誌機制，方便追蹤錯誤、事件順序與訊息內容。
    \item 將業務邏輯與各種條件判斷抽離出來，避免所有程式都塞在單一函式中。
    \item 有效管理錯誤、重試與請求頻率，以防止過載或失敗。
\end{enumerate}

若具備穩定的後端基礎架構，LINE Messaging API 能成為非常有效的數位服務入口，協助企業建立更快速、更友善的互動方式。

\chapter{本文件 LaTeX 原始碼}

以下章節展示本文件的 LaTeX 原始碼範例，包含前置設定、章節內容與程式碼排版設定。此章節會自動列入程式目錄，方便後續檢視或重複使用。

\begin{lstlisting}[language={[LaTeX]TeX}, caption={LaTeX 文件前置設定範例}, label={lst:prelude}]
\documentclass[a4paper,12pt]{report}
\usepackage[margin=2.5cm]{geometry}
\usepackage{fontspec}
\usepackage{xeCJK}
\setCJKmainfont[AutoFakeBold=2]{Microsoft JhengHei}
\usepackage{listings}
\usepackage{xcolor}
\lstset{
  numbers=left,
  breaklines=true,
  frame=single,
  backgroundcolor=\color{codebg}
}
\end{lstlisting}

\lstinputlisting[language={[LaTeX]TeX}, caption={本文件主檔 LaTeX 原始碼}, label={lst:source-main}]{main.tex}
